\documentclass[a4paper,11pt]{article}

\usepackage[a4paper, top=2.5cm, bottom=2.5cm]{geometry}
\usepackage[utf8]{inputenc}
\usepackage[french, english]{babel}

\selectlanguage{french}

\usepackage{enumitem}
\usepackage{hyperref}
\usepackage{multicol}
\usepackage{float}
\usepackage{array}
\usepackage{caption}
\usepackage{listings}

\renewcommand{\thesection}{\arabic{section}\quad---}
\renewcommand{\thesubsection}{\arabic{section}.\arabic{subsection}\quad---}

\usepackage{graphicx}

\newcolumntype{M}[1]{>{\raggedright}m{#1}}
\newcolumntype{C}[1]{>{\centering\arraybackslash}m{#1}}

%BEGIN LATEX
  \usepackage{core}

  \newcommand{\mousevox}[2]{\IfStrEq{\jobname}{\detokenize{mouse}}{#1}{#2}}

  \lstset{breaklines=true}
%END LATEX

\title{
  \blogtitle{Des ceintures tigres et panthères !}
}
\author{}
\date{}

\pagestyle{plain}
\begin{document}
\neobar{article04.fr}{nav/navi.fr}

\maketitle

\thispagestyle{first}
\begin{blogmain}
Dans l'\bloglink{article01.fr}{Article01} sur les sous-verres, nous avons appris à graver/découper au laser, à peindre au pinceau ou à la bombe aérosol.
Ce sont des pièces très bien pour commencer, pour obtenir rapidement des résultats.
Certains projets comme la fabrication de ceinture nécessitent d’apprendre à travailler les tranches et à monter des ferrures.

Je vous propose aujourd'hui de découvrir comment fabriquer des ceintures, stylisées avec des rayures de tigre ou avec un motif neutre de panthère noir.
Ce que je recherchais était une ceinture robuste, qui convienne au port d'outil comme une batterie à plomb, qui soit imperméabilisé.

Puisque je n'avais jamais conçu de patron de ceinture, j'ai étudié celle que je portais, j'ai reproduit ce plan :

\includebeltschemas{belt-35-tiger}{Tigre 3,5cm, trous traversants ronds}
\includeschemas{tab-35-tiger}{Boucle universelle simple}{1cm}

J'ai par la suite compris que les boucles que j'utilisais pouvaient être à rouleau pour mieux glisser, à ardillons doubles pour mieux supporter le poids d'outils lourds. Que leurs trous traversants devraient être de forme ovale pour mieux correspondre aux ardillons d'une boucle à rouleau, que ces trous traversants ovales pouvaient être métallisés avec des ferrures pour réduire le stress mécanique de la charge des outils.

J'ai donc pu m'adapter ce patron :

\includebeltschemas{belt-50-black-leopard}{Panthère noire 5cm, boucle à double ardillons, trous traversants ovales}
\includeschemas{ext-50-black-leopard}{Extension passant}{1.5cm}

Mes étapes de fabrications sont :
\begin{itemize}
  \item Graver le cuir avec le motif tigre
  \item Découper au laser sur la surface de cuir les contours de la ceinture 
  \item Appliquer une gomme de finition sur les tranches/la croûte, imperméabiliser le cuir
  \item Monter les ferrures
\end{itemize}

\begin{multicols}{2}

\section{Les tranches}
Les tranches sont les bords du cuir.

Mes premières versions de ceinture tigres avaient des tranches à nu, cela leurs donnes un style particulièrement brutal qui a sa beauté.
Ce n'est cependant pas une bonne idée pour la durée de vie de nos ceintures.
Ces tranches à nu sont poreuses à l'humidité, ce qui comprend notre sueur, celle-ci a des sels/acides qui sont corrosifs pour le cuir.

Un traitement recommandé est d'appliquer une gomme de finition.
Ça donnera aux tranches un aspect plus lisse qui les rendront plus étanches, et qui nous aidera à mieux les imperméabiliser.

Une gomme de finition peut aussi avoir une fonction esthétique.
Certaines gommes de finition sont colorées, d'autres peuvent être mélangés avec une teinture à base aqueuse
C'est optionnellement applicable à la croûte pour uniformiser la couleur, pour presser ses fibres.

J'ai testé la \href{https://www.decocuir.com/blogs/blog-tutoriel-patron-cuir/tokonole-gomme-cuir-brunir-tranche}{Tokonole} qui m'a convenu.

Ça s'applique avec un brunissoir.
Étant d'un naturel pressé, j'ai préféré un \href{https://www.decocuir.com/products/rouleau-depose-teinture-tranche-cuir-ceinture-anse-bords}{rouleau applicateur de peinture}.

\section{Les ferrures}
Les ferrures sont des pièces métalliques
Avec une ceinture, celles-ci sont la boucle, les rivets pour la fixer, optionnellement un passant et des œillets pour les ardillons de boucle.
Ces métaux peuvent par exemple être de l'acier, du laiton plaqué nickel, du zamak, de l'Inox, \dots

Aujourd'hui, j'utilise idéalement des ferrures en Inox 316L, elles sont plus résistantes à l'oxydation, ne nécessitant pas de placages qui se craquelleraient lors de chocs répétés.
Celles-ci nécessitant cependant des techniques de moulages plus complexes, elles sont donc simplifiées, moins répandus/plus coûteuses.
J'en ai beaucoup appris sur ce sujet grâce à cet article \href{https://la-mode-a-l-envers.loom.fr/blog/cette-ceinture-vous-survivra}{Cette ceinture vous survivra} de la société \href{https://www.loom.fr/}{Loom}

\end{multicols}

J'en ai appris d'avantage sur le travail du cuir.
Ça m'a demandé de travailler avec la longueur maximale du plateau de découpe.

\includeimage{photo-35-tiger}{Modèle Tigre}
\includeimage{photo-50-black-leopard}{Modèle Panthères}

J'ai en bonus fabriqué des portes clés de ceintures personnalisées.

De futurs projets avec le cuir pourraient impliquer de la couture, de l'imperméabilisation avec de la graisse marine.
\end{blogmain}
\end{document}
